% DRAFT results subsection for RQ6. The development part is final; the test part carries TODO-FINAL markers.

\subsection{RQ6: Guided generation with verify-and-repair}\label{sec:res-rq6}

\textbf{Development.} With the first prompt version, all twelve development loops (three models, four projects) passed the gate, after at most three repairs. Four of them nevertheless failed the independent smoke test, and the reasons shaped the second version. In one, every pod was Ready but the front end listened on port 8079 while its container port and Service pointed at port 80. Without a readiness probe, Ready said nothing about whether the process served traffic, and the documentation never states the port. This led to the listening-port check (vii) and to a rule requiring a readiness probe on every served port. In a second, the model took a dashboard version that the documentation mentions only in a history note (``Starting from \ldots v2.0.0'') as the image tag, a side effect of the rule to use documented values; every other configuration of that project, from every model, used the current release. In the remaining two, the dashboard reached its metrics scraper through the API server's service proxy, which fails for a named Service port (Section~\ref{sec:res-rq4}), or the scraper was left out altogether. Two further rules came from failures that repair had fixed: a service that mounts a ReadWriteOnce volume cannot run two replicas on different nodes, and a read-only root file system needs writable home and cache directories. Developing the loop also exposed a flaw in our own feedback: a configuration placed in \texttt{kube-system} was sent the logs of the cluster's components, including the names of other projects. We restricted feedback to objects the configuration created and re-ran the affected rounds.

With the second version, Claude's four development configurations passed the gate and the smoke test at the first attempt, where the first version had needed a repair for one and had failed the smoke test for another. The dashboard configuration now set the scraper's address explicitly, as the new rule asks. Gemini passed n8n at the first attempt, where it had needed a repair before, and the voting app as before; on the dashboard it used the current release but left the scraper out (\TODO{TODO-FINAL: Gemini Sock Shop, delayed by provider outages}). Qwen's first attempts got worse: none passed the smoke test, against two of four with the production prompt. Repair recovered one, and for two projects the provider rejected the repair request as too large. We froze the prompt at the second version, for two reasons. The remaining failures were an omitted component, which no generic cluster signal reports, and provider limits, which no prompt can change. And a third version would have re-run every development cell from the budget the held-out evaluation needed.

\textbf{Held-out projects.} \TODO{TODO-FINAL: Table~\ref{tab:rq6}; per model: smoke and Ready rates under neutral, production, guided first attempt and final attempt; McNemar tests; repairs needed; gate-passed-but-smoke-failed cases; security of final attempts (PSS, Kubescape) vs production; provider-stopped loops.}
